\section{Software integration and exact interfaces}
\label{sec:integration}

The supplied scanner changes are optional. The public default remains
\code{reduction="standard"}, which selects the existing sparse
minimum-fill cancellation. Earlier structural tests and recognition
decisions retain their existing order. Selecting a new reducer changes
the chain-complex calculation only if execution reaches that stage.
The dihedral-cover classifier is a separate topology kernel; it is not
silently substituted for a normal-surface or knot-recognition routine.

\subsection{Scanner modules and their responsibilities}

The new scanner code is organized by mathematical responsibility.
The following filenames are relative to \path{fast/fastunknot/}.

\begin{center}
\small
\begin{tabularx}{\textwidth}{@{}p{0.31\textwidth}X@{}}
\toprule
Module & Contract\\
\midrule
\code{graded.py} &
Retains raw quantum shifts, checks homogeneous degrees, predicts the
full scalar survivor profile, and supplies support-only reference
reducers.\\
\code{binary_contraction.py} &
Preserves the inherited scalar-contraction basis and its global
packed-column representation.\\
\code{scalar_components.py} &
Splits the scalar graph into components, handles isolated objects
and identity pairs directly, and returns sparse index columns in
the same survivor basis.\\
\code{corridor.py} &
Constructs the typed transfer graph, prunes irrelevant paths,
evaluates either direction, and installs the full minimal
differential through scanner-compatible classes.\\
\bottomrule
\end{tabularx}
\end{center}

The quantum array is aligned with the scanner's object slots.
Ordinary cancellation can leave dead slots, and the grading code
preserves that convention. Crossing expansion installs the new shifts
before dispatching to the chosen reducer. The shifts are the raw
scanning shifts of equation~\eqref{eq:corridor-q-recurrence}; no
diagram-wide orientation normalization is silently applied.

The direct kernel \code{corridor_transfer(scan)} requires a valid
homogeneous chain complex in the scanner category. It returns the
surviving matching, homological-degree, and quantum-shift arrays, the
full output differential, the scalar contraction, component details,
and work counters. The original object and differential arrays are
not replaced by this function, although coefficient caches can be
populated. \code{CorridorScan} performs the replacement after the
kernel completes.

\subsection{Public recognition and raw-homology options}

Two reduction names are added to \code{khovanov_rank},
\code{recognize}, and their corresponding command-line commands:
\begin{description}
\item[\code{corridor}]
Compute the full survivor differential by scalar contraction and
automatic bidirectional transfer through the retained corridors at
every reduction stage.
\item[\code{corridor-adaptive}]
Start with the existing sparse reducer and permit
\(\max(256,4N)\) candidate Schur-update pairs at the beginning of
the stage, where \(N\) is the current live object count. If the
allowance is exceeded, retain completed pivots and transfer the
partially reduced complex.
\end{description}

The allowance is a heuristic for when to pay for another
representation. It has no competitive-time guarantee. It also
differs from the older \code{adaptive} reducer: that method tries a
zero-differential survivor-profile certificate and resumes ordinary
cancellation when the certificate is inconclusive. The new adaptive
method can finish with a nonzero radical differential.

A raw-homology call bypasses the recognition front ends and therefore
provides a direct way to exercise the selected reducer:

\begin{lstlisting}[language=Python]
from fastunknot import Diagram, khovanov_rank

diagram = Diagram.from_braid(3, [1, -2] * 4)
result = khovanov_rank(
    diagram.pd,
    reduction="corridor-adaptive",
    max_objects=50000,
    seconds=30,
    check_d_squared=True,
)
print(result["rank"], result["by_degree"])
\end{lstlisting}

The returned public quantities retain their prior conventions:
\code{rank} is the unreduced \(\F\)-rank,
\code{reduced_rank} is half of it for a knot, and
\code{by_degree} uses the raw homological degrees.
The result records the nonstandard reduction name and its counters.
The direct graded scanner additionally exposes
\code{ranks_by_bidegree()}; its live-object counts are homology
dimensions only once the differential has been shown to vanish,
as on a closed minimal complex.

The corresponding command-line examples, run inside \path{fast/},
are:

\begin{lstlisting}
python -m fastunknot khovanov examples/conway.json \
  --reduction corridor --check-d2

python -m fastunknot recognize examples/conway.json \
  --reduction corridor-adaptive \
  --seconds 30 --max-objects 50000
\end{lstlisting}

The recognition command can finish in an earlier certified stage,
so a chosen reduction option does not imply that a scan occurred.
Inspect the returned method and scan evidence when measuring it.
The option \code{--no-reduction} retains its previous meaning:
it disables Reidemeister preprocessing. It is unrelated to the
\code{--reduction} chain-cancellation selector.

\subsection{Research controls for direction and scalar setup}

The public reduction names use
\code{mode="auto"}, \code{prune=True},
\code{direction_policy="reach"}, and
\code{scalar_engine="global"}. The sparse scalar variant is exposed
through the direct scanner/kernel API, allowing the additional setup
choices to remain independently measurable:

\begin{lstlisting}[language=Python]
from time import monotonic
from fastunknot.corridor import CorridorScan
from fastunknot.ordering import best_scan_order

scan = CorridorScan(
    mode="auto",
    prune=True,
    direction_policy="ports",
    scalar_engine="components",
    max_objects=50000,
    deadline=monotonic() + 30,
)
for index in best_scan_order(diagram.pd):
    scan.add_crossing(diagram.pd[index])
print(scan.total_rank())
\end{lstlisting}

This example continues with the validated nonempty diagram above.
For ordinary use, the public raw-homology API also handles its
established zero-crossing convention and other dispatch details.

The \code{mode} values are \code{forward}, \code{reverse}, and
\code{auto}. Disabling pruning supplies an ablation rather than a
different homotopy invariant. The \code{reach} policy constructs
packed endpoint sets and minimizes the finer structural propagation
bound separately in each weak component. The \code{ports} policy
uses Boolean reachability and compares the component's numbers of
input and output ports. It avoids endpoint-mask setup but can choose
a direction with a worse fine-grained bound. Neither policy is
described as an oracle for elapsed time.

The \code{global} scalar engine retains the frozen construction.
The \code{components} engine uses the sparse direct-sum construction
proved above. These controls can also be passed to
\code{AdaptiveCorridorScan}. There is no additional public
\code{corridor-compressed} reduction name in this release.

\subsection{Compatibility and resource semantics}

The supported combinations are explicit:

\begin{center}
\small
\begin{tabularx}{\textwidth}{@{}p{0.28\textwidth}X@{}}
\toprule
Option or interface & Corridor requirement or behavior\\
\midrule
Scanner backend &
The standard scanner backend; the shared, saturated, Euler, and
twist backends are outside this integration.\\
Pivot and algebra &
\code{pivot="minfill"}, \code{algebra="bits"}, and
\code{self_inverse=True} in the raw API.\\
Racing &
\code{race=1}; competing scan processes are not combined with these
reducers.\\
Coefficient composition &
The standard, component, and forced component-dense coefficient
engines remain selectable subject to their existing limits.\\
Tail contraction &
Permitted. The requested final crossings are added without
cancellation, then the closed scalar complex is finished by the
existing binary rank calculation.\\
Homological windows &
The public corridor recognition options require
\code{window_radius=None}. The separate window command retains
its existing reduction choices.\\
\bottomrule
\end{tabularx}
\end{center}

Incompatible raw-API combinations raise \code{ValueError}. The CLI
rejects incompatible backend, algebra, pivot, racing, sharing, twist,
and corridor/window combinations as option errors. These guards
prevent a requested reducer from being silently replaced by an
unsupported implementation.

The object ceiling is checked before crossing expansion allocates
the predicted number of objects. It is a ceiling on that count,
not a limit on total Python memory or on all auxiliary scalar
matrices. Time limits are cooperative: scalar setup, graph
construction, reachability, and propagation poll the ordinary
resource checks. They are not operating-system hard deadlines.

The transfer computes its replacement before installing new object
and map arrays. A resource failure during that calculation therefore
does not install an incomplete minimal differential. Adaptive
reduction retains any complete sparse pivots already performed.
This statement concerns the reducer; it does not assert that every
allocation or entire crossing-expansion operation is reversible.

The raw \code{khovanov_rank} API propagates \code{ScanLimit} and
\code{MemoryError}. The recognition pipeline catches both and
returns \code{UNKNOWN} with resource evidence. The raw CLI likewise
reports a resource failure rather than a knot verdict. Exhausting
an object or time allowance never means that a knot has been found,
and an empty set of usable shortcuts never means that a diagram is
an unknot.

\subsection{Counters and reproducible comparisons}

The \code{corridor_} counters expose original and retained graph
sizes, component counts, selected directions, structural forward
and reverse bounds, actual edge propagations, radical-edge
propagations, coefficient xors, and nontrivial compositions.
\code{corridor_switches}, when present, records adaptive transitions
from sparse cancellation. Stage histories retain the input and
survivor counts.

The structural-bound assertion checks that the counted
propagations do not exceed the chosen bound. That assertion does
not count all execution costs: the component's ordered vertices
must still be visited, and scalar/graph setup remains necessary.
The distinction is reflected in
equation~\eqref{eq:full-stage} and in the negative timings.

The package keeps both frozen predecessors, the measured source
hashes, full input descriptions and crossing orders, raw timing
batches, and exact matrix digests. The reproduction wrapper restores
the required frozen corridor in a temporary copy when replaying the
nine-arm experiment, so historical measurements can be reproduced
without overwriting the delivered implementation. The separate
compressed-setup experiment uses the frozen corridor as an explicit
baseline. The ordinary test suite and the two inexpensive support
audits can be run independently of the timing campaigns.

\subsection{The covering classifier is an explicit separate contract}

The geometric implementation is
\path{dihedral_covers/dihedral_cover.py}. Its input schema specifies
an orientable bordered surface by genus and boundary count, or a
nonorientable bordered surface by crosscap and boundary counts,
together with a positive sheet count and one affine permutation for
each canonical free generator. Shifts are reduced modulo the sheet
count. The canonical schema supplies the boundary words and
orientation character, so accepted data do not depend on checking
an omitted closed-surface relation.

\begin{lstlisting}
python dihedral_covers/dihedral_cover.py \
  dihedral_covers/examples/three_types.json
\end{lstlisting}

The Python entry points are \code{classify_cover},
\code{component_key}, and \code{boundary_lift_key}. The first returns
at most three component-family records, their multiplicities,
covering degrees, surface topology, complete boundary-degree
profiles, and cover-isomorphism type identifiers. Distinct
exceptional residue records can share a type identifier while
remaining separately queryable. The two marked-query functions
identify the component and lifted boundary circle containing a
specified fibre point without expanding the sheets.

The field \code{orientation_interval_bundle} describes the
canonical interval bundle with orientable total space. It does not
specify an arbitrary extra twisting character. Unknown fields,
incorrect generator counts, invalid signs, inappropriate surface
parameters, and booleans supplied in place of integers are rejected.

This schema certifies an abstract finite surface cover. Extracting
that schema from a normal surface, certifying an embedding in a
three-manifold, and retaining arbitrary external attaching maps
are separate tasks. No current scanner or recognizer dispatches to
this classifier merely because a large integer weight appears.
The integration point for a future compressed cutting backend must
supply the covering presentation and its attachment data
explicitly, as required in Section~\ref{sec:accounting}.
